\chapter{Etude du backend Spring Boot}

\section{Stack technique}

Le backend est un projet Spring Boot (parent \texttt{spring-boot-starter-parent} version 3.2.4) avec Java 17.

\begin{table}[H]
\centering
\begin{tabular}{p{0.35\textwidth} p{0.55\textwidth}}
\toprule
\textbf{Composant} & \textbf{Technologie} \\
\midrule
Framework principal & Spring Boot 3.2.4 \\
Langage & Java 17 \\
Web/API & spring-boot-starter-web \\
Data access & spring-boot-starter-data-jpa \\
Securite & spring-boot-starter-security + JWT (jjwt) \\
Validation & spring-boot-starter-validation \\
Temps reel & spring-boot-starter-websocket + STOMP \\
Documentation API & springdoc-openapi-starter-webmvc-ui \\
Messaging interne & spring-jms + Artemis Jakarta \\
Base de donnees & MySQL (driver 8.0.33) \\
\bottomrule
\end{tabular}
\caption{Synthese stack backend}
\end{table}

\section{Configuration serveur et donnees}

Le backend expose par defaut:
\begin{itemize}[leftmargin=*]
  \item \texttt{server.port=\$\{PORT:8091\}};
  \item \texttt{server.servlet.context-path=/football-academy}.
\end{itemize}

La datasource configuree cible MySQL avec JPA/Hibernate. Le stockage des uploads repose sur \texttt{file.storage.location=uploads} et une exposition statique des ressources.

\section{Securite et autorisations}

\subsection*{JWT et filtres}
La classe \texttt{JwtAuthenticationFilter} supporte:
\begin{itemize}[leftmargin=*]
  \item header \texttt{Authorization: Bearer <token>};
  \item cookie \texttt{accessToken};
  \item fallback parametre query (debug).
\end{itemize}

\subsection*{Regles de securite}
La \texttt{SecurityConfig} applique une politique stateless et distingue:
\begin{itemize}[leftmargin=*]
  \item chemins publics (auth, assets, uploads, websocket);
  \item routes admin restreintes a \texttt{ROLE\_ADMIN};
  \item routes \texttt{/api/**} authentifiees.
\end{itemize}

\section{Organisation fonctionnelle REST}

Les controleurs REST couvrent plusieurs domaines metier:
\begin{itemize}[leftmargin=*]
  \item authentification et profil;
  \item gestion activites/matchs/calendar;
  \item joueurs, entraineurs, parents, divisions, paiements;
  \item notifications;
  \item recherche globale;
  \item chat et contacts.
\end{itemize}

\section{Chat temps reel backend}

Le chat combine:
\begin{itemize}[leftmargin=*]
  \item endpoints REST \texttt{/api/chat/...} pour conversations/messages;
  \item WebSocket STOMP endpoint \texttt{/ws};
  \item broker simple avec prefixes \texttt{/topic} et \texttt{/queue};
  \item file JMS interne pour l'envoi asynchrone des messages.
\end{itemize}

Cette architecture favorise la fluidite de l'interface chat et limite les rafraichissements manuels.

\section{Build, tests et qualite}

Le \texttt{pom.xml} inclut:
\begin{itemize}[leftmargin=*]
  \item \texttt{maven-surefire-plugin} pour tests unitaires;
  \item \texttt{maven-failsafe-plugin} pour tests integration;
  \item \texttt{jacoco-maven-plugin} pour couverture;
  \item profil \texttt{analysis} avec SpotBugs et Checkstyle.
\end{itemize}

\section{Conteneurisation et CI/CD}

\subsection*{Docker}
Le \texttt{Dockerfile} utilise un build multi-stage:
\begin{enumerate}[leftmargin=*]
  \item stage Maven pour compiler le jar;
  \item stage runtime Temurin pour executer \texttt{app.jar}.
\end{enumerate}

\subsection*{Docker Compose}
\texttt{docker-compose.yml} definit le service app, mappe le port \texttt{8091:8091} et monte le volume uploads.

\subsection*{GitLab CI/CD}
Le pipeline \texttt{.gitlab-ci.yml} execute:
\begin{itemize}[leftmargin=*]
  \item stage \texttt{test};
  \item stage \texttt{build};
  \item stage \texttt{docker-build} avec push registry.
\end{itemize}

\section{Point d'attention DB}

Le document \texttt{README-ONEDRIVE-DB.md} rappelle qu'un lien OneDrive ne remplace pas un serveur DB. Une base relationnelle reelle (MySQL/PostgreSQL) reste obligatoire pour l'execution applicative.
